% ===========================================================================
% preambulo-uis.tex
% Formato de trabajo de grado UIS, APA 7.
% Ver AGENTS.md, secciones 1 a 9. No cambiar valores de forma sin instruccion.
% ===========================================================================

% ---------- Codificacion, idioma y fuente ----------
\usepackage[utf8]{inputenc}
\usepackage[T1]{fontenc}
% es-tabla: "Tabla" en vez de "Cuadro".
% es-nosectiondot: sin punto tras numeros de seccion (se repone solo en nivel 1).
\usepackage[spanish,es-tabla,es-nosectiondot]{babel}
\usepackage{amsmath}               % antes de newtxmath
\usepackage{newtxtext}            % Times New Roman
\usepackage{newtxmath}            % incluye los simbolos de amssymb
\usepackage{csquotes}

% ---------- Pagina: carta, margenes 2,54 cm uniformes ----------
\usepackage[letterpaper,margin=2.54cm,headheight=15pt]{geometry}

% ---------- Interlineado doble, sangria 1,27 cm, justificado ----------
\usepackage{setspace}
\doublespacing
\setlength{\parindent}{1.27cm}
\setlength{\parskip}{0pt}

% ---------- Paquetes de contenido ----------
\usepackage{graphicx}
\graphicspath{{figuras/}}
\usepackage{booktabs}
\usepackage{array}
\usepackage{makecell}
\usepackage[table]{xcolor}
\usepackage{enumitem}
\setlist{nosep}
\usepackage{tikz}
\usetikzlibrary{arrows.meta,positioning,fit,backgrounds,babel}

\definecolor{morado1}{HTML}{7B6BAF}
\definecolor{azul2}{HTML}{A4C2F4}
\definecolor{verde1}{HTML}{76A5AF}
\definecolor{lima1}{HTML}{B6D7A8}

% ---------- Cornisa y paginacion ----------
% Cornisa: izquierda, MAYUSCULA SOSTENIDA, maximo 50 caracteres.
% Paginacion arabiga desde la portadilla (pagina 1).
\usepackage{fancyhdr}
\newcommand{\cornisa}{OPEN VOCABULARY CHANGE DETECTION} % 32 caracteres
\pagestyle{fancy}
\fancyhf{}
\fancyhead[L]{\cornisa}
\fancyhead[R]{\thepage}
\renewcommand{\headrulewidth}{0pt}
\fancypagestyle{plain}{%
  \fancyhf{}%
  \fancyhead[L]{\cornisa}%
  \fancyhead[R]{\thepage}%
  \renewcommand{\headrulewidth}{0pt}}

% ---------- Jerarquia de titulos (5 niveles APA) ----------
\usepackage{titlesec}
\setcounter{secnumdepth}{5}
% Nivel 1: centrado, negrita, punto tras el numero.
\titleformat{\section}[block]{\centering\bfseries\normalsize}{\thesection.}{0.5em}{}
\titlespacing*{\section}{0pt}{2\baselineskip}{2\baselineskip}
% Nivel 2: margen izquierdo, negrita.
\titleformat{\subsection}[block]{\bfseries\normalsize}{\thesubsection}{0.5em}{}
\titlespacing*{\subsection}{0pt}{\baselineskip}{\baselineskip}
% Nivel 3: margen izquierdo, negrita y cursiva.
\titleformat{\subsubsection}[block]{\bfseries\itshape\normalsize}{\thesubsubsection}{0.5em}{}
\titlespacing*{\subsubsection}{0pt}{\baselineskip}{\baselineskip}
% Nivel 4: con sangria, negrita, termina en punto, texto en la misma linea.
\titleformat{\paragraph}[runin]{\bfseries\normalsize}{\theparagraph}{0.5em}{}[.]
\titlespacing*{\paragraph}{1.27cm}{\baselineskip}{0.5em}
% Nivel 5: con sangria, negrita y cursiva, termina en punto.
\titleformat{\subparagraph}[runin]{\bfseries\itshape\normalsize}{\thesubparagraph}{0.5em}{}[.]
\titlespacing*{\subparagraph}{1.27cm}{\baselineskip}{0.5em}

% Seccion sin numero que inicia en hoja nueva y entra a la Tabla de Contenido
% (Introduccion, Referencias Bibliograficas, Apendices).
\newcommand{\seccionsinnumero}[1]{%
  \clearpage
  \phantomsection
  \section*{#1}%
  \addcontentsline{toc}{section}{#1}}

% ---------- Figuras y tablas: leyenda ARRIBA ----------
\usepackage{caption}
\captionsetup{
  position=top,
  labelsep=newline,
  labelfont=bf,
  textfont=it,
  justification=raggedright,
  singlelinecheck=false,
  font=doublespacing,
  skip=6pt}
% Nota bajo figura o tabla: "Nota." en cursiva.
\newcommand{\nota}[1]{\par\vspace{4pt}{\raggedright\textit{Nota.} #1\par}}

% Nota al pie con llamada fija (* y ** en Resumen y Abstract).
\newcommand{\notapie}[2]{{\renewcommand{\thefootnote}{#1}\footnote{#2}}}

% Marcador visible para contenido que el autor debe completar.
\newcommand{\pendiente}[1]{\textcolor{red}{\textbf{[PENDIENTE:} #1\textbf{]}}}

% ---------- Tabla de contenido y listas ----------
\usepackage{tocloft}
\addto\captionsspanish{%
  \renewcommand{\contentsname}{Tabla de Contenido}%
  \renewcommand{\listfigurename}{Lista de Figuras}%
  \renewcommand{\listtablename}{Lista de Tablas}%
  \renewcommand{\tablename}{Tabla}%
  \renewcommand{\figurename}{Figura}}
\renewcommand{\cfttoctitlefont}{\hfill\bfseries\normalsize}
\renewcommand{\cftaftertoctitle}{\hfill}
\renewcommand{\cftloftitlefont}{\hfill\bfseries\normalsize}
\renewcommand{\cftafterloftitle}{\hfill}
\renewcommand{\cftlottitlefont}{\hfill\bfseries\normalsize}
\renewcommand{\cftafterlottitle}{\hfill}
\renewcommand{\cftsecaftersnum}{.}          % solo el nivel 1 lleva punto
\renewcommand{\cftsecfont}{\normalfont}
\renewcommand{\cftsecpagefont}{\normalfont}
\renewcommand{\cftsecleader}{\cftdotfill{\cftdotsep}}
\renewcommand{\cftfigpresnum}{Figura }
\renewcommand{\cfttabpresnum}{Tabla }
\setlength{\cftfignumwidth}{5.5em}
\setlength{\cfttabnumwidth}{5em}
\setlength{\cftbeforesecskip}{0pt}

% ---------- Referencias: APA 7 con biblatex ----------
\usepackage[style=apa,backend=biber]{biblatex}
\DeclareLanguageMapping{spanish}{spanish-apa}
\addbibresource{referencias.bib}
% "y" en lugar de "&" en citas y lista (spanish-apa.lbx deja smartand=0).
\AtBeginDocument{%
  \setcounter{smartand}{2}%
  \DeclareDelimFormat[parencite,textcite,bib,biblist]{finalnamedelim}{\addspace\bibstring{and}\space}%
}
\setlength{\bibhang}{1.27cm}   % sangria francesa
\setlength{\bibitemsep}{0pt}   % el doble espacio lo da \doublespacing

% ---------- hyperref al final ----------
\usepackage[hidelinks]{hyperref}
